\section{Задание}

Моделируем информационный центр. В информационный центр приходят клиенты (пользователи) через интервал времени 10 ± 2 минуты. Если все три имеющихся оператора заняты, клиенту отказывают в обслуживании. Операторы имеют разную производительность и могут обеспечивать обслуживание среднего запроса от пользователя за 20 ± 5, 40 ± 10 и 40 ± 20 ед. времени (минут). Клиенты стараются занять свободного оператора с максимальной производительностью. Полученные запросы сдаются в накопитель, откуда выбираются на обработку. На первый компьютер — от первого и второго операторов, на второй — от третьего. Время обработки запроса в компьютерах — 15 и 30 минут соответственно. Смоделировать процесс обработки 300 запросов. Определить вероятность отказа.

\section{Теоретическая часть}

\subsection{Равномерное распределение}

Случайная величина $X$ имеет равномерное распределение на отрезке $[a, b]$, если её плотность распределения $f(x)$ равна:

\begin{equation}
p(x) = \begin{cases}
\frac{1}{b-a}, & \text{если } a \leq x \leq b \\
0, & \text{иначе}
\end{cases}
\end{equation}

При этом функция распределения $F(x)$ равна:

\begin{equation}
F(x) = \begin{cases}
0, & x < a \\
\frac{x-a}{b-a}, & a \leq x \leq b \\
1, & x > b
\end{cases}
\end{equation}

Обозначение: $X \sim R[a, b]$.

Интервал времени между появлением $i$-й и $(i-1)$-й заявки по равномерному закону вычисляется следующим образом:

\begin{equation}
T_i = a + (b - a) \cdot R
\end{equation}

где $R$ — псевдослучайное число из интервала $[0, 1]$.

В GPSS для команды \texttt{GENERATE} с параметрами $A$ и $B$ (например, \texttt{GENERATE 10,2}) интервал генерируется по формуле $T = A \pm B$, что эквивалентно равномерному распределению на интервале $[A-B, A+B]$.

\subsection{Переменные и уравнение имитационной модели}

\textbf{Эндогенные переменные:}

\begin{itemize}
    \item время обработки задания $i$-ым оператором;
    \item время решения задания на $j$-ом компьютере.
\end{itemize}

\textbf{Экзогенные переменные:}

\begin{itemize}
    \item $n_0$ — число обслуженных клиентов;
    \item $n_1$ — число клиентов, получивших отказ.
\end{itemize}

Вероятность отказа в обслуживании клиента вычисляется как:

\begin{equation}
P = \frac{n_1}{n_0 + n_1}
\end{equation}

\subsection{GPSS}

Язык GPSS — общецелевая система моделирования.

Транзакты представляют собой описание динамических процессов в реальных системах. Они могут описывать как реальные физические объекты, так и нефизические. Транзакты можно генерировать и уничтожать в процессе моделирования.

Динамическими объектами являются транзакты, которые представляют собой единицы исследуемых потоков и производят ряд определённых действий, продвигаясь по фиксированной структуре.

Операционный объект. Блоки задают логику функционирования системы и определяют маршрут движения транзактов между объектами аппаратной категории.

К статическим объектам относятся очереди и таблицы, служащие для оценок влияющих характеристик.

Рассмотрим основные команды:

\begin{enumerate}
    \item \textbf{GENERATE} — команда, вводящая транзакты в модель
    \item \textbf{TERMINATE} — команда, удаляющая транзакт
    \item \textbf{QUEUE} — команда, помещающая транзакт в конец очереди
    \item \textbf{DEPART} — команда, удаляющая транзакт из очереди
    \item \textbf{SEIZE} — команда, занимающая канал обслуживания
    \item \textbf{RELEASE} — команда, освобождающая канал обслуживания
    \item \textbf{ADVANCE} — команда, задерживающая транзакт
    \item \textbf{TRANSFER} — команда, изменяющая движение транзакта в модели
    \item \textbf{GATE} — вспомогательный блок, проверяющий состояния устройств. \texttt{GATE NU} — устройство не занято
    \item \textbf{SAVEVALUE} — команда, сохраняющая значение в переменную
    \item \textbf{START} — команда, управляющая процессом моделирования
\end{enumerate}

\section{Описание реализации}

\subsection{Структура программы}

В листинге 1 представлена реализация информационного центра на языке GPSS.

\includelistingpretty
    {lab7.gps}
    {gpss}
    {Реализация информационного центра на языке GPSS}

Программа реализует следующую логику:

\begin{enumerate}
    \item Генерация клиентов с интервалом $10 \pm 2$ минуты (лимит 300 клиентов)
    \item Каскадная проверка доступности операторов с приоритетом (от самого быстрого к самому медленному):
    \begin{itemize}
        \item Если Оператор 1 свободен — обслуживание $20 \pm 5$ минут, запрос на Компьютер 1
        \item Иначе, если Оператор 2 свободен — обслуживание $40 \pm 10$ минут, запрос на Компьютер 1
        \item Иначе, если Оператор 3 свободен — обслуживание $40 \pm 20$ минут, запрос на Компьютер 2
        \item Иначе — отказ в обслуживании (метка \texttt{fail})
    \end{itemize}
    \item Обработка запросов на компьютерах:
    \begin{itemize}
        \item Компьютер 1: очередь \texttt{queue\_1}, ��ремя обработки 15 минут
        \item Компьютер 2: очередь \texttt{queue\_2}, время обработки 30 минут
    \end{itemize}
    \item Подсчёт статистики:
    \begin{itemize}
        \item Сохранение количества отказов в переменную \texttt{nfail}
        \item Вычисление вероятности отказа по формуле: \texttt{prob = nfail / (success + nfail)}
    \end{itemize}
\end{enumerate}

Ключевая особенность реализации — использование блоков \texttt{GATE NU} для последовательной проверки доступности операторов и реализации приоритетного выбора.

\section{Практическая часть}

В листинге 2 представлен результат работы программы.

\includelistingpretty
    {output.txt}
    {text}
    {Результат выполнения программы}

Основные результаты моделирования:

\begin{itemize}
    \item \textbf{Время моделирования:} 6980.179 единиц модельного времени
    \item \textbf{Всего клиентов:} 300
    \item \textbf{Успешно обслужено:} 231 клиент (77.0\%)
    \item \textbf{Получили отказ:} 69 клиентов (23.0\%)
    \item \textbf{Вероятность отказа:} 0.230 (23.0\%)
\end{itemize}

\textbf{Распределение нагрузки по операторам:}

\begin{itemize}
    \item \textbf{Оператор 1} (самый быстрый): обслужил 121 клиента (40.3\%), утилизация 34.5\%
    \item \textbf{Оператор 2}: обслужил 59 клиентов (19.7\%), утилизация 33.8\%
    \item \textbf{Оператор 3} (самый медленный): обслужил 51 клиента (17.0\%), утилизация 31.2\%
\end{itemize}

\textbf{Загрузка компьютеров:}

\begin{itemize}
    \item \textbf{Компьютер 1}: 180 запросов, утилизация 38.7\%, максимальная очередь 32
    \item \textbf{Компьютер 2}: 51 запрос, утилизация 21.9\%, максимальная очередь 115
\end{itemize}

\textbf{Анализ результатов:}

Вероятность отказа составила \textbf{23.0\%}, что означает, что почти каждый четвёртый клиент не получил обслуживания из-за занятости всех операторов.

Наблюдается значительный дисбаланс в системе:

\begin{enumerate}
    \item Операторы загружены примерно на треть времени (~33\%), но при этом 23\% клиентов получают отказ. Это связано с отсутствием очередей к операторам — при одновременной занятости всех трёх операторов клиент мгновенно получает отказ.

    \item Критическая ситуация с Компьютером 2: максимальная очередь достигла 115 запросов, среднее время ожидания составило 1722.8 минут (более суток), при этом утилизация компьютера всего 21.9\%. Это объясняется тем, что Оператор 3 обслуживает медленно (40±20 минут) и неравномерно, создавая «пачки» запросов.

    \item Оператор 1, как самый быстрый, обслужил больше всего клиентов (121), что подтверждает корректность реализации приоритетного выбора.
\end{enumerate}

\section{Вывод}

В ходе выполнения лабораторной работы была разработана модель информационного центра на языке GPSS с тремя операторами разной производительности и двумя компьютерами для обработки запросов.

Программа корректно реализует заданную систему: клиенты выбирают свободного оператора с максимальной производительностью, запросы маршрутизируются на соответствующие компьютеры, ведётся подсчёт отказов и вычисляется вероятность отказа.

Определена вероятность отказа в обслуживании — \textbf{23.0\%}. Это высокий показатель, свидетельствующий о недостаточной пропускной способности системы. Выявлены узкие места системы:

\begin{itemize}
    \item Отсутствие очередей к операторам приводит к высокой вероятности отказов при относительно низкой загрузке
    \item Компьютер 2 работает в режиме критической перегрузки очереди при низкой утилизации
    \item Дисбаланс между скоростью работы операторов и компьютеров
\end{itemize}

Для улучшения характеристик системы рекомендуется добавить очереди к операторам, увеличить количество операторов или оптимизировать распределение запросов между компьютерами.
